\section{Solution Novelty Analysis}
\label{sec:idea_novelty}

The preceding evaluations establish how well agents optimize and what shapes their performance, but a high score does not reveal whether an agent discovered a new idea or assembled established techniques.
Recent studies raise the same concern: research agents tend to remain close to prior work or recombine existing methods, while genuinely original ideas remain rare even when search explicitly targets diversity and novelty~\citep{tang2026narrow,antoniades2026heuresis}.
To characterize what current auto research agents actually produce, we analyze the best of three solutions for every model--task pair, yielding $252$ solutions.
For each solution, we extract the initial-to-final code diff, commit history, and experiment journal, and use Claude-Opus-4.8 with a fixed rubric provided in Appendix~\ref{app:novelty_judge_rubric} to classify it into one of the eight mutually exclusive categories shown in Figure~\ref{fig:idea_novelty}.
To minimize false positives in the central \emph{novel-approach} category, we manually review every candidate and retain the label only when the core idea clearly goes beyond established approaches for the task.
Accordingly, our conclusions about novelty primarily characterize the AI-for-AI optimization setting and may not extend to more open-ended scientific discovery.

\begin{figure}[t]
\centering
\includegraphics[width=\linewidth]{Figures/idea_novelty.pdf}
\caption{Novelty analysis of 252 best-of-three solutions. Left: distribution across the eight solution categories after Opus-4.8 classification. Right: the three novel approaches retained after manual review.}
\label{fig:idea_novelty}
\end{figure}

\textbf{Agents improve artifacts primarily by composing established techniques, while genuine novelty is rare.}
Composition-stacking, which layers multiple established algorithmic and engineering optimizations onto a standard approach, is the largest category for every model and accounts for $111$ of $252$ solutions ($44.0\%$).
By comparison, novel approaches are rare: after manual review, only three solutions ($1.2\%$) retain this label.
More strikingly, $16$ solutions ($6.3\%$) exploit evaluation-specific shortcuts, more than five times the novel count, with GPT-5.5 accounting for eight of these cases.
Thus, when agents depart from standard techniques, they are more likely to exploit loopholes in the evaluation protocol than to produce a validated novel approach.

\textbf{Novel approaches do not concentrate in the highest-performing models and arise through task-specific reframing rather than new technical primitives.}
One might naturally expect higher-performing models such as Opus-4.7 and GPT-5.5 to produce more novel approaches, yet the three validated cases come from GLM-5.2, Kimi-K2.7-Code, and LongCat-2.0, which occupy different positions in the overall ranking.
Specifically, GLM constructs an ancilla-free comparator by combining Fredkin-based split-and-restore with algebraic normal form, Kimi reframes next-frame prediction around optical flow and residual warping, and LongCat identifies a small set of BatchNorm bits that acts as an architectural chokepoint.
In each case, novelty lies not in inventing a new technical primitive, but in identifying a task-specific insight and using familiar components in a way that standard approaches do not suggest.
